% part: intuitionistic-logic
% chapter: propositions-as-types
% section: normalization

\documentclass[../../../include/open-logic-section]{subfiles}

\begin{document}

\olfileid[id]{int}{pty}{nor}

\olsection{Normalisasi}

Dalam bagian ini kita membuktikan bahwa, melalui suatu urutan reduksi,
setiap deduksi dapat direduksi menjadi deduksi normal; hasil ini disebut
\emph{sifat normalisasi}. Kita akan menggunakan korespondensi proposisi
sebagai tipe: kita menunjukkan bahwa setiap term bukti dapat direduksi
menjadi bentuk normal; normalisasi bagi !!{derivation}s deduksi alami
kemudian mengikuti.

Pertama-tama kita mendefinisikan beberapa fungsi yang mengukur
kompleksitas term. \emph{Panjang}~$\len{!A}$ suatu !!a{formula}s didefinisikan
oleh
\begin{align*}
  \len{p} &= 0\\
  \len{!A \land !B} &= \len{!A} + \len{!B} + 1\\
  \len{!A \lor !B} &= \len{!A} + \len{!B} + 1 \\
  \len{!A \lif !B} &= \len{!A} + \len{!B} + 1.
\end{align*}
Kompleksitas suatu redex~$M$ diukur oleh \emph{peringkat
potongannya}~$\cutrank{M}$:
\begin{align*}
  \cutrank{(\lambd[\typeof{x}{!A}][\typeof{N}{!B}])Q} &=
  \len{!A} + \len{!B} + 1 \\
  \cutrank{\ande{i}{\andi{\typeof{M}{!A}}{\typeof{N}{!B}}}} &=
  \len{!A} + \len{!B} + 1 \\
  \cutrank{\ore{\ori{i}{!A_{i}}{\typeof{M}{!A_i}}}
    {\typeof{x_1}{!A_1}}{\typeof{N_1}{!C}}
    {\typeof{x_2}{!A_2}}{\typeof{N_2}{!C}}} &=
  \len{!A_1} + \len{!A_2} + 1.
\end{align*}
Kompleksitas suatu term bukti diukur oleh redex paling kompleks di
dalamnya, dan bernilai $0$ jika term itu normal:
\[
  \maxrank{M} = \max \bigl(\{\cutrank{N} \mid
    N \text{ adalah subterm dari } M
    \text{ dan merupakan redex}\} \cup \{0\}\bigr).
\]

\begin{lem}\ollabel{lem:subst}
  Jika $\Subst{M}{\typeof{N}{!A}}{\typeof{x}{!A}}$ merupakan redex dan
  $M \not\ident x$, maka salah satu kasus berikut berlaku:
  \begin{enumerate}
  \item $M$ sendiri merupakan redex; atau
  \item $M$ berbentuk $\ande{i}{x}$, dan $N$ berbentuk
    $\andi{P_1}{P_2}$;
  \item $M$ berbentuk $\ore{x}{x_1}{P_1}{x_2}{P_2}$, dan $N$ berbentuk
    $\ori{i}{}{Q}$; atau
  \item $M$ berbentuk $x Q$, dan $N$ berbentuk $\lambd[x][P]$.
  \end{enumerate}
  Dalam kasus pertama, $\cutrank{\Subst{M}{N}{x}} = \cutrank{M}$;
  dalam kasus-kasus lain,
  $\cutrank{\Subst{M}{N}{x}} = \len{!A}$.
\end{lem}
\begin{proof}
  Bukti melalui induksi pada~$M$.
  \begin{enumerate}
  \item Jika $M$ merupakan variabel tunggal $y$ dan $y \not\ident x$,
    maka $\Subst{y}{N}{x}$ ialah $y$, sehingga bukan redex.
  \item Jika $M$ berbentuk $\andi{N_1}{N_2}$, atau $\lambd[x][N]$, atau
    $\ori{i}{!A}{N}$, maka
    $\Subst{M}{\typeof{N}{!A}}{\typeof{x}{!A}}$ juga berbentuk demikian,
    sehingga bukan redex.
  \item Jika $M$ berbentuk $\ande{i}{P}$, kita mempertimbangkan dua kasus.
    \begin{enumerate}
      \item Jika $P$ berbentuk $\andi{P_1}{P_2}$, maka $M \ident
        \ande{i}{\andi{P_1}{P_2}}$ merupakan redex, dan jelas
        \[
        \Subst{M}{N}{x} \ident
        \ande{i}{\andi{\Subst{P_1}{N}{x}}{\Subst{P_2}{N}{x}}}
        \]
        juga merupakan redex. Peringkat potongannya sama.
      \item Jika $P$ merupakan variabel tunggal, variabel itu haruslah
        $x$ agar substitusi menghasilkan redex, dan $N$ harus berbentuk
        $\andi{P_1}{P_2}$. Sekarang tinjau
        \[
        \Subst{M}{N}{x} \ident
        \Subst{\ande{i}{x}}{\andi{P_1}{P_2}}{x},
        \]
        yang sama dengan $\ande{i}{\andi{P_1}{P_2}}$. Peringkat
        potongannya sama dengan $\len{!A}$.
    \end{enumerate}
  \end{enumerate}
  Kasus $\ore{N}{x_1}{N_1}{x_2}{N_2}$ dan $P Q$ serupa.
\end{proof}

\begin{lem}
  Jika $M$ berkontraksi menjadi $M'$, dan $\cutrank{M} >
  \cutrank{N}$ untuk setiap subterm-redex sejati $N$ dari $M$, maka
  $\cutrank{M} > \maxrank{M'}$.
\end{lem}

\begin{proof}
  Bukti dengan membedakan kasus.
  \begin{enumerate}
  \item Jika $M$ berbentuk $\ande{i}{\andi{M_1}{M_2}}$, maka $M'$ ialah
    $M_i$; karena setiap subterm dari $M_i$ juga merupakan subterm sejati
    dari $M$, klaim tersebut berlaku.
  \item Jika $M$ berbentuk
    $(\lambd[\typeof{x}{!A}][N])\typeof{Q}{!A}$, maka $M'$ ialah
    $\Subst{N}{\typeof{Q}{!A}}{\typeof{x}{!A}}$. Tinjau suatu redex dalam
    $M'$. Redex itu bersesuaian dengan suatu redex dalam $N$ yang
    mempunyai peringkat potongan sama, yang menurut asumsi lebih kecil
    daripada $\cutrank{M}$; atau peringkat potongannya sama dengan
    $\len{!A}$, yang menurut definisi lebih kecil daripada
    $\cutrank{(\lambd[x^{!A}][N])Q}$.
  \item Jika $M$ berbentuk
    \[
    \ore{\ori{i}{}{\typeof{N}{!A_i}}}
        {\typeof{x_1}{!A_1}}{\typeof{N_1}{!C}}
        {\typeof{x_2}{!A_2}}{\typeof{N_2}{!C}},
    \]
    maka $M' \ident \Subst{N_i}{N}{\typeof{x_i}{!A_i}}$. Tinjau suatu
    redex dalam~$M'$. Redex itu bersesuaian dengan suatu redex dalam
    $N_i$ yang mempunyai peringkat potongan sama, yang menurut asumsi
    lebih kecil daripada $\cutrank{M}$; atau peringkat potongannya sama
    dengan $\len{!A_i}$, yang menurut definisi lebih kecil daripada
    $\cutrank{\ore{\ori{i}{}{\typeof{N}{!A_i}}}
      {\typeof{x_1}{!A_1}}{\typeof{N_1}{!C}}
      {\typeof{x_2}{!A_2}}{\typeof{N_2}{!C}}}$.
  \end{enumerate}
\end{proof}

\begin{thm}
  Semua term bukti tereduksi menjadi bentuk normal; semua
  !!{derivation}s tereduksi menjadi !!{derivation}s normal.
\end{thm}

\begin{proof}
  Klaim kedua mengikuti klaim pertama. Kita membuktikan klaim pertama
  melalui induksi lengkap pada $m=\maxrank{M}$, dengan $M$ suatu term
  bukti.
  \begin{enumerate}

  \item Jika $m=0$, $M$ sudah normal.
  \item
    Jika tidak, kita melakukan induksi pada~$n$, yaitu banyaknya redex
    dalam $M$ yang mempunyai peringkat potongan sama dengan~$m$.
    \begin{enumerate}
    \item Jika $n=1$, pilih redex $N$ sedemikian rupa sehingga
      $m = \cutrank{N} > \cutrank{P}$ bagi setiap subterm sejati $P$
      yang juga merupakan redex. Redex semacam itu harus ada karena
      setiap term hanya mempunyai berhingga banyak subterm.

      Misalkan $N'$ adalah hasil reduksi $N$. Menurut lema,
      $\maxrank{N'} < \maxrank{N}$, sehingga $n$, yaitu banyaknya redex
      dengan $\cutrank=m$, berkurang. Dengan demikian $m$ berkurang
      (sebesar $1$ atau lebih), dan kita dapat menerapkan hipotesis
      induksi untuk~$m$.
    \item Untuk langkah induksi, andaikan $n>1$. Prosesnya serupa, tetapi
      $n$ hanya berkurang menjadi suatu bilangan positif sehingga $m$
      tidak berubah. Kita cukup menerapkan hipotesis induksi untuk~$n$.
    \end{enumerate}
  \end{enumerate}
\end{proof}

Fakta bahwa term-$\lambd$ bertipe ternormalisasi berarti bahwa setiap
term semacam itu \emph{memiliki bentuk normal}. Jika kita memandang
reduksi-$\beta$ term-$\lambd$ menuju bentuk normalnya sebagai pelaksanaan
komputasi, dan bentuk normal sebagai nilai komputasinya, ini berarti
bahwa setiap komputasi dalam kalkulus-$\lambd$ bertipe akhirnya berhenti
dan menghasilkan suatu nilai. Lebih lanjut, preservasi tipe memastikan
bahwa nilai tersebut memiliki tipe yang tepat. Misalnya, jika $N$
bertipe $!A \lif !B$ (yakni, ia merupakan fungsi yang seharusnya bekerja
pada argumen bertipe~$!A$ dan menghasilkan nilai bertipe~$!B$), dan kita
menerapkannya kepada term~$M$ bertipe~$!A$ (masukan), maka $(NM)$
tereduksi menjadi term~$N'$ (keluaran), dan keluaran tersebut dijamin
bertipe~$!B$.

Sebenarnya, term dalam kalkulus-$\lambd$ bertipe bukan hanya tereduksi
menjadi bentuk normal menurut satu cara penerapan reduksi-$\beta$ kepada
subterm; \emph{setiap} cara menerapkan reduksi-$\beta$ kepada subterm
akhirnya berhenti dalam bentuk normal. Sifat ini disebut
\emph{normalisasi kuat}. Artinya, bukan hanya ada suatu cara untuk
menjamin bahwa komputasi akhirnya berhenti dan menghasilkan nilai,
melainkan \emph{setiap} cara menjalankan komputasi pada akhirnya
menghasilkan nilai.

Bagaimana kita mengetahui bahwa cara-cara berbeda menjalankan komputasi
menghasilkan nilai yang \emph{sama}? Sejauh ini kita hanya mengetahui
(dari preservasi tipe) bahwa semua nilai yang dihasilkan bertipe sama.
Kalkulus-$\lambd$ bertipe memiliki sifat penting lain: ia
\emph{konfluen}, atau mempunyai sifat \emph{Church--Rosser}. Jika
$N \red N_1$ dan $N \red N_2$, maka ada suatu term~$N'$ sedemikian rupa
sehingga $N_1 \red N'$ dan $N_2 \red N'$. Ingat bahwa $N \red N'$
berarti $N'$ dihasilkan dari nol atau lebih langkah reduksi~$\redone$.
Jika $N_1$ berada dalam bentuk normal, satu-satunya term $N'$ sedemikian
rupa sehingga $N_1 \red N'$ ialah $N_1$ sendiri (melalui nol langkah
reduksi~$\redone$). Karena itu, jika $N_1$ dan $N_2$ keduanya berada
dalam bentuk normal, maka $N_1 = N' = N_2$. Dengan kata lain, karena
$\red$ ternormalisasi kuat, setiap cara mereduksi suatu term akhirnya
menghasilkan bentuk normal. Karena $\red$ bersifat Church--Rosser, setiap
dua bentuk normal sama. Jadi, komputasi dalam kalkulus-$\lambd$ bertipe
selalu berhenti dan nilainya (bentuk normalnya) sama tanpa bergantung pada
cara komputasi berlangsung.

Membuktikan bahwa suatu relasi bersifat konfluen biasanya sulit. Namun,
kita dapat dengan mudah menetapkan syarat yang lebih lemah berikut:

\begin{prop}[Sifat Church--Rosser Lemah]
Jika $N \redone N_1$ dan $N \redone N_2$, maka ada term~$N'$ sedemikian
rupa sehingga $N_1 \red N'$ dan $N_2 \red N'$.
\end{prop}

\begin{proof}
  Kedua term $N_1$ dan $N_2$ masing-masing dihasilkan dari $N$ dengan
  mengganti redex $M_1$ atau $M_2$ dengan hasil reduksinya. Redex
  $M_1$ dan $M_2$ adalah subterm dari $N$ dan tidak dapat bertumpang
  tindih sebagian. Jadi, salah satunya merupakan subterm dari yang lain,
  atau keduanya muncul di tempat berbeda yang tidak bertumpang tindih
  dalam~$N$. Tinjau kasus kedua: $N$ berbentuk $N(M_1,M_2)$. Mengganti
  $M_1$ dengan hasil reduksinya $M_1'$ menghasilkan $N_1$, yaitu
  $N(M_1',M_2)$, sedangkan mengganti $M_2$ dengan hasil reduksinya $M_2'$
  menghasilkan $N_2$, yaitu $N(M_1,M_2')$. Keduanya tereduksi menjadi
  $N(M_1',M_2')$. Untuk kasus lainnya, andaikan $M_2$ merupakan subterm
  dari~$M_1$ (kasus sebaliknya simetris), dan bedakan kasus menurut jenis
  redex~$M_1$. Misalnya, jika $M_1 \ident
  \proj{1}{\pair{O_1}{O_2}}$, ia tereduksi menjadi $O_1$. Jika $M_2$
  merupakan subterm $O_1$, mengonversi $M_2$ menghasilkan $O_1'$.
  Jadi, mengonversi $M_1$ terlebih dahulu lalu $M_2$ menghasilkan
  $O_1'$. Namun, mengonversi $M_2$ terlebih dahulu menghasilkan
  $\proj{1}{\pair{O_1'}{O_2}}$, dan term ini juga tereduksi
  menjadi~$O_1'$.
\end{proof}

\begin{lem}[Lema Newman]
Jika $\red$ ternormalisasi kuat dan memiliki sifat Church--Rosser lemah,
maka $\red$ bersifat konfluen.
\end{lem}

\begin{proof}
  Karena $\red$ ternormalisasi kuat, setiap barisan reduksi maksimal
  berakhir dalam bentuk normal. Bentuk normal unik jika dan hanya jika
  $\red$ bersifat konfluen. Andaikan suatu term $N$ mempunyai dua bentuk
  normal berbeda, $N'$ dan $N''$, yang dihasilkan oleh barisan reduksi
  \begin{align*}
  N & = N_1' \redone N_2' \redone \dots \redone N_k' = N'
    \text{ dan}\\
  N & = N_1'' \redone N_2'' \redone \dots \redone N_l'' = N''.
  \end{align*}
  (Kedua barisan reduksi harus mempunyai panjang $>0$, sebab jika panjang
  salah satunya nol maka $N$ sudah berada dalam bentuk normal dan
  tidak dapat mempunyai bentuk normal lain.)

  Jika $N_2'=N_2''$, term ini mempunyai dua bentuk normal berbeda,
  $N'$ dan $N''$. Jika $N_2' \neq N_2''$, sifat Church--Rosser lemah
  memberikan suatu $N'''$ sedemikian rupa sehingga $N_2' \red N'''$
  dan $N_2'' \red N'''$. Karena normalisasi kuat, $N'''$ tereduksi
  menjadi suatu bentuk normal~$Q$. Jika $Q \neq N'$, maka $N_2'$
  mempunyai dua bentuk normal berbeda, $N'$ dan $Q$; jika $Q=N'$,
  maka, karena $N'\neq N''$, $N_2''$ mempunyai dua bentuk normal
  berbeda, $N''$ dan $Q$. Jadi, apabila $N$ mempunyai dua bentuk normal
  berbeda, terdapat $N_1$ sedemikian rupa sehingga $N \redone N_1$ dan
  $N_1$ mempunyai dua bentuk normal berbeda. Jelas kita dapat meneruskan
  cara ini dan menghasilkan $N \redone N_1 \redone N_2 \redone \dots$,
  tetapi hal itu mustahil karena $\redone$ ternormalisasi kuat.
\end{proof}

\end{document}
